\documentclass[12pt, a4paper]{ctexart}
\usepackage{fontspec}
\usepackage{xcolor}
\usepackage{hyperref}
\usepackage{geometry}
\usepackage{enumitem}
\usepackage{parskip}
\usepackage{titlesec}
\usepackage{tcolorbox}
\tcbuselibrary{breakable}
\tcbset{autoparskip}

\usepackage{setspace}
\usepackage{listings}
\usepackage{amssymb}

\geometry{left=2.5cm, right=2.5cm, top=2.5cm, bottom=2.5cm}

\setmainfont{Times New Roman}
\setCJKmainfont{SimSun}

\hypersetup{
    colorlinks=true,
    linkcolor=blue!70!black,
    urlcolor=cyan,
    pdftitle={泛型-逐题精讲解义},
    pdfauthor={专升本-fifine老师}
}

\definecolor{myblue}{RGB}{30,100,200}
\definecolor{lightblue}{RGB}{230,240,255}
\definecolor{darkblue}{RGB}{0,50,120}

\newtcolorbox{bluebox}[1][]{
    breakable,
    colback=lightblue,
    colframe=myblue,
    coltitle=darkblue,
    fonttitle=\bfseries,
    title={#1},
    boxrule=0.8pt,
    arc=4pt,
    left=6pt,
    right=6pt,
    top=6pt,
    bottom=6pt,
    before skip=8pt,
    after skip=8pt
}

\usepackage{etoolbox}
\BeforeBeginEnvironment{bluebox}{\par\vfill}%
\AfterEndEnvironment{bluebox}{\par\vfill}%

\titleformat{\section}
  {\normalfont\Large\bfseries\color{myblue}}{\thesection}{1em}{}
\titlespacing*{\section}{0pt}{3.5ex plus 1ex minus .2ex}{1.5ex plus .2ex}

\title{\textcolor{myblue}{\textbf{泛型 - 逐题精讲解义}}}
\author{专升本-fifine老师 教学组}
\date{2026年03月05日}

\lstset{
    language=Java,
    basicstyle=\ttfamily\small,
    keywordstyle=\color{blue},
    commentstyle=\color{green!60!black},
    stringstyle=\color{orange},
    numbers=left,
    numberstyle=\tiny\color{gray},
    frame=single,
    breaklines=true,
    columns=fullflexible,
    showstringspaces=false
}

\newcommand{\code}[1]{\texttt{#1}}

\begin{document}

\maketitle
\thispagestyle{empty}
\newpage

\section{泛型}

\begin{enumerate}[label=\textbf{\arabic*.}]

	\item \textbf{以下关于Java中的泛型的说法错误的是}
	      \begin{enumerate}[label=\Alph*.]
		      \item 可以提高代码的安全性
		      \item 可以减少类型转换
		      \item 泛型类型在运行时会被擦除
		      \item 泛型可以用于基本数据类型
	      \end{enumerate}

	      \begin{bluebox}{答案与深度解析}
		      \textbf{答案：D. 泛型可以用于基本数据类型}

		      \textbf{【核心认知框架】}
		      \begin{itemize}[label={}, leftmargin=*, nosep]
			      \item \textbf{题目本质}：考查泛型的基本特性和限制
			      \item \textbf{关键区分点}：泛型只支持引用类型，不支持基本类型
			      \item \textbf{命题人意图}：测试对泛型类型擦除和类型参数限制的理解
		      \end{itemize}

		      \textbf{【选项原理深度拆解】}

		      \begin{itemize}[leftmargin=*, nosep, label={}]
			      \item[A] \textbf{可以提高代码的安全性 — 正确}
			            \begin{itemize}[label={}, leftmargin=*, nosep]
				            \item \textbf{编译时类型检查}：泛型在编译时进行类型验证，避免运行时ClassCastException
				            \item \textbf{代码示例}：
				                  \begin{lstlisting}[language=Java]
// 不使用泛型（不安全）
List list = new ArrayList();
list.add("Hello");
String s = (String) list.get(0);  // 需要强制转换
// list.add(100);  // 如果添加Integer，运行时抛异常

// 使用泛型（安全）
List<String> list2 = new ArrayList<>();
list2.add("Hello");
String s2 = list2.get(0);  // 无需转换
// list2.add(100);  // 编译错误！类型不安全
				      \end{lstlisting}
				            \item \textbf{字节码层面}：编译器插入checkcast指令确保类型安全
			            \end{itemize}

			      \item[B] \textbf{可以减少类型转换 — 正确}
			            \begin{itemize}[label={}, leftmargin=*, nosep]
				            \item \textbf{自动类型转换}：从泛型集合中获取元素时自动插入checkcast
				            \item \textbf{代码对比}：
				                  \begin{lstlisting}[language=Java]
// 无泛型：手动转换
List rawList = new ArrayList();
rawList.add(100);
Integer num = (Integer) rawList.get(0);

// 有泛型：自动处理（编译器插入checkcast）
List<Integer> genericList = new ArrayList<>();
genericList.add(100);
Integer num2 = genericList.get(0);
// 编译后字节码中自动包含checkcast指令
				      \end{lstlisting}
			            \end{itemize}

			      \item[C] \textbf{泛型类型在运行时会被擦除 — 正确}
			            \begin{itemize}[label={}, leftmargin=*, nosep]
				            \item \textbf{类型擦除机制}：泛型类型信息在编译后被擦除，运行时无法获取
				            \item \textbf{JLS规范}：JLS §4.6 类型擦除
				            \item \textbf{字节码证据}：
				                  \begin{lstlisting}[language=Java]
// 源码
List<String> list1 = new ArrayList<>();
List<Integer> list2 = new ArrayList<>();

// 字节码（javap -c）
// 反编译后：
List list1 = new ArrayList();
List list2 = new ArrayList();
// 泛型类型信息消失，都是原始类型List

// Class对象相同
System.out.println(list1.getClass() == list2.getClass());
// 输出：true（都是ArrayList.class）
				      \end{lstlisting}
				            \item \textbf{擦除规则}：
				                  \begin{itemize}[label={}, nosep]
					                  \item 无界<T>擦除为Object
					                  \item 有界<T extends Number>擦除为Number
					                  \item 方法泛型参数也被擦除
				                  \end{itemize}
			            \end{itemize}

			      \item[D] \textbf{泛型可以用于基本数据类型 — 错误（本题答案）}
			            \begin{itemize}[label={}, leftmargin=*, nosep]
				            \item \textbf{错误原因}：泛型类型参数必须是引用类型，不能是基本类型
				            \item \textbf{JLS规范}：JLS §4.5 类型参数必须是引用类型
				            \item \textbf{字节码层面}：JVM指令集对泛型处理要求引用类型
				            \item \textbf{代码反证}：
				                  \begin{lstlisting}[language=Java]
// 编译错误！
List<int> list = new ArrayList<>();
// error: 泛型类型参数不能是基本数据类型

// 错误！
Map<int, String> map = new HashMap<>();

// 错误！
Pair<int, double> pair = new Pair<>();
				      \end{lstlisting}
				            \item \textbf{正确做法}：使用对应的包装类
				                  \begin{lstlisting}[language=Java]
// 使用包装类 Integer
List<Integer> list = new ArrayList<>();
list.add(100);  // 自动装箱 int -> Integer
int num = list.get(0);  // 自动拆箱 Integer -> int

// 其他基本类型的包装类
// int -> Integer
// double -> Double
// boolean -> Boolean
// char -> Character
// byte -> Byte
// short -> Short
// long -> Long
// float -> Float
				      \end{lstlisting}
				            \item \textbf{为什么不能支持基本类型}：
				                  \begin{itemize}[label={}, nosep]
					                  \item 类型擦除后基本类型信息丢失
					                  \item JVM泛型实现在字节码层面依赖引用类型
					                  \item 避免复杂的重载和类型系统问题
				                  \end{itemize}
			            \end{itemize}
		      \end{itemize}

		      \textbf{【应背诵知识点】}
		      \begin{itemize}[label={}, nosep]
			      \item 泛型类型参数必须是引用类型
			      \item 基本类型需使用包装类（int->Integer等）
			      \item 泛型类型在运行时被擦除
			      \item 泛型提供编译时类型安全，减少强制转换
			      \item 口诀：泛型只支持引用型，基本类型要装箱，运行类型全擦除，编译安全运行无
		      \end{itemize}

		      \textbf{【命题趋势与实战策略】}
		      \textbf{考场秒杀}：看到"泛型基本数据类型"→错，必须用包装类
	      \end{bluebox}

	      \vspace{1em}

	\item \textbf{以下哪个是正确的Java泛型声明？}
	      \begin{enumerate}[label=\Alph*.]
		      \item class MyClass<T extends Number>
		      \item class MyClass<T super Number>
		      \item class MyClass<T implements Number>
		      \item class MyClass<T>
	      \end{enumerate}

	      \begin{bluebox}{答案与深度解析}
		      \textbf{答案：A. class MyClass<T extends Number>}

		      \textbf{【核心认知框架】}
		      \begin{itemize}[label={}, leftmargin=*, nosep]
			      \item \textbf{题目本质}：考查泛型类型参数的限定语法
			      \item \textbf{关键区分点}：类型上界可以用extends，不能用super或implements
			      \item \textbf{命题人意图}：测试对泛型类型边界语法的准确掌握
		      \end{itemize}

		      \textbf{【选项原理深度拆解】}

		      \begin{itemize}[leftmargin=*, nosep, label={}]
			      \item[A] \textbf{class MyClass<T extends Number> — 正确}
			            \begin{itemize}[label={}, leftmargin=*, nosep]
				            \item \textbf{语法规则}：<T extends 上界类型>
				            \item \textbf{适用场景}：限定类型参数必须是指定类型或其子类
				            \item \textbf{代码示例}：
				                  \begin{lstlisting}[language=Java]
// T必须是Number或其子类（Integer, Double等）
class Calculator<T extends Number> {
    public double add(T a, T b) {
        return a.doubleValue() + b.doubleValue();
    }
}

// 正确使用
Calculator<Integer> intCalc = new Calculator<>();
Calculator<Double> doubleCalc = new Calculator<>();

// 编译错误！String不是Number的子类
Calculator<String> stringCalc = new Calculator<>();
// error: 类型参数String不在类型变量T的范围内
				      \end{lstlisting}
				            \item \textbf{多边界}：
				                  \begin{lstlisting}[language=Java]
// &连接多个边界（类在前，接口在后）
class MyClass<T extends Number & Comparable<T>> {
    // T必须同时是Number子类和Comparable接口的实现类
}

// Integer既是Number子类，又实现了Comparable
Calculator<Integer> c = new Calculator<>();  // 正确
				      \end{lstlisting}
				            \item \textbf{仅接口边界}：
				                  \begin{lstlisting}[language=Java]
// 限定T必须实现Comparable接口
class Sorter<T extends Comparable<T>> {
    public void sort(T[] array) {
        Arrays.sort(array);  // 使用Comparable的compareTo
    }
}
// 即使边界是接口，也用extends关键字
				      \end{lstlisting}
			            \end{itemize}

			      \item[B] \textbf{class MyClass<T super Number> — 错误}
			            \begin{itemize}[label={}, leftmargin=*, nosep]
				            \item \textbf{错误原因}：泛型类型参数声明不能用super限定
				            \item \textbf{说明}：super用于通配符下界，不用于类型参数声明
				            \item \textbf{代码对比}：
				                  \begin{lstlisting}[language=Java]
// 类型参数声明不能用super
class MyClass<T super Number> {}  // 编译错误

// 通配符可以使用super
void method(List<? super Number> list) {
    // ? super Number表示Number或其父类
    // ? extends Number表示Number或其子类
}
// super只能在通配符中使用，不能在类型参数声明中使用
				      \end{lstlisting}
			            \end{itemize}

			      \item[C] \textbf{class MyClass<T implements Number> — 错误}
			            \begin{itemize}[label={}, leftmargin=*, nosep]
				            \item \textbf{错误原因}：无论限界是类还是接口，都用extends
				            \item \textbf{JLS规范}：类型参数边界统一使用extends
				            \item \textbf{代码反证}：
				                  \begin{lstlisting}[language=Java]
// 编译错误！类型参数限界不能用implements
class MyClass<T implements Number> {}
// error: 期望是类型，而不是接口

// 正确：无论限界是类还是接口都用extends
class MyClass<T extends Number> {}  // Number是类
class MyClass<T extends Comparable<T>> {}  // Comparable是接口
// implements仅用于类实现接口的声明处
				      \end{lstlisting}
			            \end{itemize}

			      \item[D] \textbf{class MyClass<T> — 语法正确但不是本题最佳答案}
			            \begin{itemize}[label={}, leftmargin=*, nosep]
				            \item \textbf{说明}：无界泛型声明，没有类型限制
				            \item \textbf{代码示例}：
				                  \begin{lstlisting}[language=Java]
// 无界泛型
class Box<T> {
    private T value;
    public void set(T value) { this.value = value; }
    public T get() { return value; }
}

// 可以接受任何引用类型
Box<String> stringBox = new Box<>();
Box<Integer> intBox = new Box<>();
Box<Object> objectBox = new Box<>();
				      \end{lstlisting}
			            \end{itemize}
		      \end{itemize}

		      \textbf{【通配符vs类型参数限界】}
		      \begin{itemize}[label={}, nosep]
			      \item 类型参数：<T extends Number>在泛型类/方法声明中
			      \item 通配符：<? extends Number>在变量声明、方法参数中
			      \item 下界：<? super Number>只能用于通配符，不能用于类型参数
		      \end{itemize}

		      \textbf{【应背诵知识点】}
		      \begin{itemize}[label={}, nosep]
			      \item 类型参数限界：<T extends 上界>
			      \item 无论限界是类还是接口都用extends
			      \item 不能用super或implements
			      \item 多边界用&连接：& (类 & 接口 & 接口)
			      \item 通配符下界：<? super 类型>（调用时使用）
			      \item 口诀：类型限界extends，类和接口都用它，super只能通配符用，多界连接用符号&
		      \end{itemize}

		      \textbf{【命题趋势与实战策略】}
		      \textbf{考场秒杀}：看到泛型类型参数限界→用extends（无论类还是接口）
		      \textbf{记忆}：implements只在类实现接口时，类型限界统一用extends
	      \end{bluebox}

\end{enumerate}

\end{document}
